% !TeX root = ../教材题库_全部条目_转写工作稿.tex
% Source: XJ-B1-C06-RV-P265-Y1455
% 第6章 / 章末复习 复习题六 / 教材 p.265、266 / PDF 第272、273页
% OCR未校对：公式与文字均可能有误；未完成数学TeX转写；禁止直接发给学生
% 原题及相邻原解截图（核对用）：../题目原图/XJ-B1-C06-RV-P265-Y1455-01.webp, ../题目原图/XJ-B1-C06-RV-P265-Y1455-02.webp
% 本文件只在转写完成后作为题干使用；OCR 条目尚不是可靠题干。
\begin{bankitem}{复习题六 第9题}{第6章 / 章末复习 复习题六 / 教材 p.265、266 / PDF 第272、273页}{XJ-B1-C06-RV-P265-Y1455}{OCR未校对：公式与文字均可能有误；未完成数学TeX转写；禁止直接发给学生}
\begin{ocrdraft}
[教材p.265 / PDF第272页]
\par
9．甲、乙、丙、丁四人参加全省数学竞赛，初赛时四人分别处于第95，第
\par
96，
\par
第97，第98百分位数处，已知参加初赛的总人数为50000，并且取成绩前
\par
1500名进入决赛，试问：上述四人中有哪几位可参加决赛？
\par
[教材p.266 / PDF第273页]
\par
温故而知新
\end{ocrdraft}
\end{bankitem}
